\section{Conclusions}
\label{sec:conclusions}

This project looked at how HGCAL clustering behaves when photons no
longer point back to the interaction point, and which specific
reconstruction features cost performance. Along the way I improved the displaced-particle gun, added displacement and timing observables
to the CMSSW validation, made the association thresholds configurable, and
produced the two-dimensional performance maps used throughout. Comparing
shower-axis estimators gave an approximately \SI{10}{cm} resolution on the
back-extrapolated displacement of the most energetic Trackster. Small fragments do
much worse than this, so it should not be read as a number that describes every
reconstructed object.

A large part of the work went into explaining why default \cluethreed{} loses
single displaced photons. Because it projects everything from the origin, two points
on the same shower axis look transversely separated by roughly $R|\Delta z|/z$,
and this apparent separation grows with displacement. The loss comes in two
stages: around $R=\SI{200}{cm}$ the energy carried by fragments whose
nearest-higher link skips two layers rises and then levels off, which produces
the intermediate plateau, and at larger displacement one-layer breaks take over
and drive the final decline. Widening the seed distance from
\SI{1.8}{cm} to \SI{2.1}{cm}, with the density radius fixed, pushes the
fragmentation out to larger displacement and recovers substantially more of the
shower.

The two-photon studies make the case that energy recovery cannot be judged on
its own. For photons \SI{4}{cm} apart, the energy linked across the two showers
keeps rising with displacement even where the object-based merge rate is
falling, so fragmentation can make the merge rate look better while the mixing
underneath is getting worse. In PU200 the injected signal photon behaves much
as it does without pileup. The inclusive truth population is a different story:
its particle mix, kinematics and shower sizes all shift with resolved
displacement, and charged particles bend in the magnetic field, so its
displacement curve cannot be read as a clean scan of production radius.

Comparing metrics through \cluestering{} showed both promise and limits.
Several Cartesian configurations stay efficient on single displaced photons
across the whole scan, including in pileup, but the Euclidean ones merge the
\SI{4}{cm} photon pair heavily, so high efficiency alone does not mean the two
particles were actually resolved. The gnomonic configuration was the
most encouraging, giving higher efficiency than default \cluethreed{} at large
displacement while keeping a low merge rate on the two-photon test, though this
still needs further studies and its EE-only
coverage is why its inclusive pileup efficiency looks so low.

The natural next step is to tune coordinates, metric weights, distance scales
and density thresholds together, judging them against energy recovery,
fragmentation and cross-shower mixing at once. It also needs to be tested on more
realistic inputs: the controlled photons here are injected just before HGCAL,
so they say nothing about upstream interactions or a real LLP lifetime. The
generation and validation tools built for this study give a starting point for
that work, and the geometric picture points to the specific clustering
decisions worth changing.
